\section{Method}

\subsection{Finite-Support Aggregate Likelihood}

Let a bag contain instances $x_1,\ldots,x_n$.
For each instance, a neural predictor emits a categorical distribution over latent integer contributions:
\begin{equation}
  d_i(k) = p_\theta(z_i = k \mid x_i), \qquad k \in \{a,\ldots,b\}.
\end{equation}
The observed bag label is an aggregate
\begin{equation}
  Y = \sum_{i=1}^n z_i.
\end{equation}
The exact bag likelihood is the coefficient of $Y$ in the convolution of atomic PMFs:
\begin{equation}
  p_\theta(Y=y \mid x_{1:n})
  =
  \left(d_1 * d_2 * \cdots * d_n\right)(y).
\end{equation}
Training minimizes the negative log likelihood
\begin{equation}
  \mathcal{L}_{\mathrm{agg}}(\theta)
  =
  -\log p_\theta(Y=y \mid x_{1:n}).
\end{equation}
The support of the aggregate PMF is finite and contiguous:
\begin{equation}
  \mathrm{supp}(Y) = \{na, na+1,\ldots, nb\}.
\end{equation}
All experiments use exact likelihoods over this support; approximations enter only through the neural instance model, not through sampling latent assignments.

\subsection{Training Objective}

For each minibatch, the instance model first produces atomic PMFs for every instance in every bag.
The aggregate layer then deterministically maps these atoms to an exact PMF over the observed bag label.
The training loss is the negative log probability assigned to the observed aggregate,
\begin{equation}
  \mathcal{L}(\theta)
  =
  -\frac{1}{B}\sum_{b=1}^B
  \log
  \left(d_{b1} * \cdots * d_{bn_b}\right)(y_b).
\end{equation}
Gradients are backpropagated through the convolutional aggregate layer into the instance network.
Instance labels are never used for training in the weak-supervision experiments; they are used only for evaluation of hidden instance recovery.
This is different from expected-value regression baselines, which regress the scalar expectation $\sum_i \mathbb{E}[z_i]$ to the aggregate label and therefore discard the higher-order distributional information in the count or sum.

\subsection{Special Cases}

\paragraph{Binary counts.}
For Bernoulli target indicators,
\begin{equation}
  d_i = (1-p_i, p_i), \qquad Y=\sum_i z_i.
\end{equation}
Equivalently, each instance contributes a length-two one-dimensional kernel over support $\{0,1\}$:
\begin{equation}
  \kappa_i = [1-p_i,\; p_i].
\end{equation}
The probability of observing count $y$ is the coefficient of $t^y$ in
\begin{equation}
  \prod_{i=1}^n \left((1-p_i) + p_i t\right).
\end{equation}
This recovers the standard binary count likelihood.

\paragraph{Signed counts.}
For known signs $s_i \in \{-1,+1\}$ and latent target indicators $z_i \in \{0,1\}$,
\begin{equation}
  Y = \sum_i s_i z_i.
\end{equation}
Each atom has support $\{-1,0,+1\}$ after incorporating the sign:
\begin{equation}
  d_i =
  \begin{cases}
    (1-p_i)\delta_0 + p_i\delta_{+1}, & s_i=+1,\\
    (1-p_i)\delta_0 + p_i\delta_{-1}, & s_i=-1.
  \end{cases}
\end{equation}
In array order $[-1,0,+1]$, the corresponding signed kernels are
\begin{equation}
  \kappa_i^{(+)} = [0,\;1-p_i,\;p_i],
  \qquad
  \kappa_i^{(-)} = [p_i,\;1-p_i,\;0].
\end{equation}
This setting exposes cancellation: many distinct latent assignments can yield the same observed aggregate, especially when $Y=0$.

\paragraph{Ordinal sums.}
For hidden ordinal values $z_i \in \{0,\ldots,K-1\}$, the observed label can be a scalar sum such as a digit sum.
Unlike expected-sum regression, the likelihood uses the full induced distribution over possible sums.
If a model predicts digit probabilities $p_\theta(z_i=k\mid x_i)$, the likelihood compares the observed sum against the full convolved PMF rather than only the mean $\sum_i \mathbb{E}[z_i]$.
The one-dimensional kernel for instance $i$ is the full class-probability vector
\begin{equation}
  \kappa_i = [p_{i0}, p_{i1}, \ldots, p_{i,K-1}],
\end{equation}
so the aggregate support has width $n(K-1)+1$.
For MNIST digit-sum, $K=10$ and the observed bag label is the sum of hidden digit values.

\paragraph{Histogram labels and one-vs-rest counts.}
For multiclass histograms, each class count can be treated as a one-vs-rest Bernoulli count.
Let $p_{ic}$ be the predicted probability of class $c$ for instance $i$, and let $h_c$ be the observed count of class $c$ in the bag.
For every class $c$, we form Bernoulli count kernels
\begin{equation}
  \kappa_{ic} = [1-p_{ic},\;p_{ic}]
\end{equation}
and compute a classwise count likelihood.
The one-vs-rest count objective is
\begin{equation}
  \mathcal{L}_{\mathrm{hist}}(\theta)
  =
  -\frac{1}{C}\sum_{c=1}^C
  \log
  \left[
    \prod_{i=1}^n \left((1-p_{ic}) + p_{ic}t\right)
  \right]_{t^{h_c}},
\end{equation}
where $[\cdot]_{t^{h_c}}$ denotes coefficient extraction.
This yields a PVC-style objective closely related to LLP-PVC, while sharing the same finite-support convolution backend.

\subsection{Posterior Marginals}

For binary counts, we compute count-conditioned posterior marginals
\begin{equation}
  q_i = p_\theta(z_i=1 \mid \sum_j z_j = y, x_{1:n}),
\end{equation}
using prefix/suffix count distributions.
Let $F_{i-1}(k)$ be the probability that the first $i-1$ instances contain $k$ positives, and let $B_{i+1}(k)$ be the probability that instances $i+1,\ldots,n$ contain $k$ positives.
Then
\begin{equation}
  q_i
  =
  \frac{
    p_i \sum_k F_{i-1}(k) B_{i+1}(y-1-k)
  }{
    \sum_k F_{i-1}(k) \left((1-p_i)B_{i+1}(y-k) + p_iB_{i+1}(y-1-k)\right)
  }.
\end{equation}
These marginals are used as additional information extracted from the count label after likelihood training.
They support instance recovery diagnostics, uncertainty analysis, and optional pseudo-label inspection, but they are not part of the core training objective in this paper.

\subsection{Computational Backends}

All of the above objectives are one-dimensional convolutions over integer aggregate support.
If $m^{(i)}$ denotes the aggregate PMF after processing the first $i$ atoms, then
\begin{equation}
  m^{(i)}(s)
  =
  \sum_{r=a}^{b} m^{(i-1)}(s-r)d_i(r),
  \qquad
  m^{(0)}=\delta_0.
\end{equation}
This recurrence is the same probability law whether implemented as dynamic programming, shifted dense tensor accumulation, grouped \texttt{conv1d}, or FFT polynomial multiplication.
The implementation uses several exact backends:
\begin{itemize}
  \item Sequential finite-support convolution accumulates shifted copies of the current PMF and is used for moderate binary, signed, and ordinal-sum experiments.
  \item Dynamic programming is used as a Shukla-style Bernoulli count baseline; it computes the same binary count coefficients as the convolutional view.
  \item Grouped GPU convolution packs signed Bernoulli groups into channels and applies grouped \texttt{conv1d}.
  \item Balanced FFT-tree convolution multiplies the atomic polynomials pairwise and is used for large-bag runtime experiments and multiclass one-vs-rest histogram/PVC likelihoods.
\end{itemize}

\paragraph{Grouped signed \texttt{conv1d}.}
The literal GPU \texttt{conv1d} backend is used for grouped signed binary aggregates.
Each active group starts as a point mass at zero on support $[-K_{\max},K_{\max}]$.
For one signed instance, the backend applies a length-three channelwise kernel and padding one.
Because PyTorch's \texttt{conv1d} operator is cross-correlation rather than mathematical convolution, the code stores the kernel in the order
\begin{equation}
  [p_i\mathbf{1}\{s_i=+1\},\;1-p_i,\;p_i\mathbf{1}\{s_i=-1\}],
\end{equation}
so the left kernel entry shifts probability mass toward larger aggregate values and the right entry shifts mass toward smaller values.
This is only an implementation orientation; in mathematical support order $[-1,0,+1]$ the signed kernels are the atoms shown above.

\paragraph{Multiclass and histogram implementation.}
For digit-sum, the atom width is the number of ordinal labels, e.g. ten for MNIST digits.
For histogram/PVC experiments, we avoid a full $C$-dimensional joint histogram over classes and instead compute $C$ one-vs-rest one-dimensional count likelihoods in parallel.
This keeps the aggregate support one-dimensional for each class and makes large bags practical on GPU.
The FFT-tree backend is especially useful for this setting because the same pairwise polynomial products are batched over classes and bags.
Padding or masked positions are represented by deterministic zero-contribution atoms, such as $[1,0]$ for Bernoulli counts, $[0,1,0]$ for signed support $[-1,0,+1]$, and $[1,0,\ldots,0]$ for ordinal sums.

\begin{table}[t]
\centering
\small
\caption{Exact aggregate backends.  $n$ is bag size, $K$ is atomic support width, $S=n(K-1)+1$ is aggregate support width, and $C$ is number of one-vs-rest classes.}
\label{tab:backends}
\begin{tabular}{p{0.18\linewidth}p{0.24\linewidth}p{0.20\linewidth}p{0.25\linewidth}}
\toprule
Backend & Setting & Work per bag & Use in this paper \\
\midrule
DP & Bernoulli count & $O(nS)$ & Shukla-style count baseline \\
Sequential convolution & finite support & $O(nKS)$ & main moderate-bag experiments \\
Grouped convolution & signed Bernoulli & $O(nS)$, batched groups & signed GPU implementation \\
FFT-tree & finite support / OVR & $O(CS\log S\log n)$ & large-bag runtime, PVC-style histograms \\
\bottomrule
\end{tabular}
\end{table}
